\chapter{1986年全国卷（文）}

\section{单选题}

\begin{problem}
在下列各数中，已表示成三角形式的复数是（\quad）

\choices
    {\(2\left(\cos \frac{\pi}{4} - \i\sin \frac{\pi}{4}\right)\)}
    {\(2\left(\cos \frac{\pi}{4} + \i\sin \frac{\pi}{4}\right)\)}
    {\(2\left(\sin \frac{\pi}{4} + \i\cos \frac{\pi}{4}\right)\)}
    {\(-2\left(\sin \frac{\pi}{4} - \i\cos \frac{\pi}{4}\right)\)}
\end{problem}

\begin{answer}
B
\end{answer}

\begin{problem}
函数 \(y = 5^x + 1\) 的反函数是（\quad）

\choices
    {\(y = \log_5(x + 1)\)}
    {\(y = \log_x 5 + 1\)}
    {\(y = \log_5(x - 1)\)}
    {\(y = \log_{\left(x - 1\right)} 5\)}
\end{problem}

\begin{answer}
C
\end{answer}

\begin{problem}
已知全集 \(I=\{1,2,3,4,5,6,7,8\}\)，\(A=\{3,4,5\}\)，\(B=\{1,3,6\}\)，那么集合 \(\{2,7,8\}\) 是（\quad）

\choices
    {\(A\cup B\)}
    {\(A\cap B\)}
    {\(\overline{A}\cup\overline{B}\)}
    {\(\overline{A}\cap\overline{B}\)}
\end{problem}

\begin{answer}
D
\end{answer}

\begin{problem}
函数 \(y=\sqrt{2}\sin 2x\cos 2x\) 是（\quad）

\choices
    {周期为 \(\frac{\pi}{2}\) 的奇函数}
    {周期为 \(\frac{\pi}{2}\) 的偶函数}
    {周期为 \(\frac{\pi}{4}\) 的奇函数}
    {周期为 \(\frac{\pi}{4}\) 的偶函数}
\end{problem}

\begin{answer}
A
\end{answer}

\begin{problem}
已知 \(c>0\)，在下列不等式中成立的一个是（\quad）

\choices
    {\(c>2^c\)}
    {\(c>\left(\frac{1}{2}\right)^c\)}
    {\(2^c<\left(\frac{1}{2}\right)^c\)}
    {\(2^c>\left(\frac{1}{2}\right)^c\)}
\end{problem}

\begin{answer}
D
\end{answer}

\begin{problem}
有以下 \(20\) 个数：\(87\)，\(91\)，\(94\)，\(88\)，\(93\)，\(91\)，\(89\)，\(87\)，\(92\)，\(86\)，\(90\)，\(92\)，\(88\)，\(90\)，\(91\)，\(86\)，\(89\)，\(92\)，\(95\)，\(88\)，它们的和是（\quad）

\choices
    {\(1789\)}
    {\(1799\)}
    {\(1879\)}
    {\(1899\)}
\end{problem}

\begin{answer}
B
\end{answer}

\begin{problem}
已知某正方体对角线长为 \(a\)，那么这个正方体的全面积是（\quad）

\choices
    {\(2\sqrt{2}a^2\)}
    {\(2a^2\)}
    {\(2\sqrt{3}a^2\)}
    {\(3\sqrt{2}a^2\)}
\end{problem}

\begin{answer}
B
\end{answer}

\begin{problem}
如果方程 \(x^2 + y^2 + Dx + Ey + F = 0\)（\(D^2 + E^2 - 4F > 0\)）所表示的曲线关于直线 \(y = x\) 对称，那么必有（\quad）

\choices
    {\(D = E\)}
    {\(D = F\)}
    {\(E = F\)}
    {\(D = E = F\)}
\end{problem}

\begin{answer}
A
\end{answer}

\begin{problem}
设甲是乙的充分条件，乙是丙的充要条件，丙是丁的必要条件，那么丁是甲的（\quad）

\choices
    {充分条件}
    {必要条件}
    {充要条件}
    {既不充分也不必要的条件}
\end{problem}

\begin{answer}
D
\end{answer}

\begin{problem}
在下列各图中，\(y = ax^2 + bx\) 与 \(y = ax + b\)（\(ab \neq 0\)）的图象只可能是（\quad）

\choices
    {\choicebitmap[width=0.15\paperwidth]{img/1986/national_liberal/q10_choiceA.png}}
    {\choicebitmap[width=0.15\paperwidth]{img/1986/national_liberal/q10_choiceB.png}}
    {\choicebitmap[width=0.15\paperwidth]{img/1986/national_liberal/q10_choiceC.png}}
    {\choicebitmap[width=0.15\paperwidth]{img/1986/national_liberal/q10_choiceD.png}}
\end{problem}

\begin{answer}
C
\end{answer}

\section{解答题}

\begin{problem}
求方程 \(\sqrt{25^{x^2 + x - 0.5}} = \sqrt[4]{5}\) 的解.
\end{problem}

\begin{solution}
由于 \(25=5^2\)，且 \(25^{x^2+x-\frac{1}{2}}>0\)，所以原方程等价于 \(5^{x^2+x-\frac{1}{2}}=5^{\frac{1}{4}}\). 因为底数 \(5>1\)，故 \(x^2+x-\frac{1}{2}=\frac{1}{4}\)，即 \(x^2+x-\frac{3}{4}=0\). 因式分解得 \(\left(x-\frac{1}{2}\right)\left(x+\frac{3}{2}\right)=0\)，所以原方程的解为 \(x=\frac{1}{2}\) 或 \(x=-\frac{3}{2}\).
\end{solution}

\begin{problem}
已知 \(\omega = \frac{-1 - \sqrt{3}\i}{2}\)，求 \(\omega^2 + \omega + 1\) 的值.
\end{problem}

\begin{solution}
由已知 \(\omega=\frac{-1-\sqrt{3}\i}{2}\)，得 \(\omega^2=\frac{\left(-1-\sqrt{3}\i\right)^2}{4}=\frac{1+2\sqrt{3}\i-3}{4}=\frac{-1+\sqrt{3}\i}{2}\). 因而 \(\omega^2+\omega+1=\frac{-1+\sqrt{3}\i}{2}+\frac{-1-\sqrt{3}\i}{2}+1=0\).
\end{solution}

\begin{problem}
在 \(xOy\) 平面上，\(\triangle ABC\) 的三个顶点坐标依次为 \((0,0)\)，\((1,0)\)，\((0,3)\)，将这个三角形绕 \(x\) 轴旋转一周，求所得到的几何体的体积.
\end{problem}

\begin{solution}
三角形内部由 \(x\) 轴、\(y\) 轴以及直线 \(y=3-3x\) 围成，其中直线在 \(x=0\) 时的纵坐标为 \(3\)，在 \(x=1\) 时的纵坐标为 \(0\). 绕 \(x\) 轴旋转后，所得几何体是底面半径为 \(3\)、高为 \(1\) 的圆锥，因此其体积为 \(V=\frac{1}{3}\pi\cdot3^2\cdot1=3\pi\).
\end{solution}

\begin{problem}
求 \(\displaystyle\lim_{n\to\infty}\frac{2n^2 + n + 7}{5n^2 + 4}\).
\end{problem}

\begin{solution}
分子、分母同除以 \(n^2\)，得

\[
\displaystyle\lim_{n\to\infty}\frac{2n^2+n+7}{5n^2+4}=\displaystyle\lim_{n\to\infty}\frac{2+\frac{1}{n}+\frac{7}{n^2}}{5+\frac{4}{n^2}}=\frac{2}{5},
\]

因为 \(\displaystyle\lim_{n\to\infty}\frac{1}{n}=0\) 且 \(\displaystyle\lim_{n\to\infty}\frac{1}{n^2}=0\).
\end{solution}

\begin{problem}
求 \(\left(2x^3 - \frac{1}{x^2}\right)^5\) 展开式中的常数项.
\end{problem}

\begin{solution}
设展开式中取了 \(j\) 个 \(-\frac{1}{x^2}\) 因式，其中 \(j=0,1,\ldots,5\)，则相应项为

\[
\mathrm C_5^j(2x^3)^{5-j}\left(-\frac{1}{x^2}\right)^j=\mathrm C_5^j2^{5-j}(-1)^j x^{15-5j}.
\]

要得到常数项，必须使 \(x\) 的幂为零，即 \(15-5j=0\)，所以 \(j=3\). 此时常数项的系数为 \(\mathrm C_5^3 2^2(-1)^3=-40\)，故展开式中的常数项为 \(-40\).
\end{solution}

\begin{problem}
求与椭圆 \(\frac{x^2}{9} + \frac{y^2}{4} = 1\) 有公共焦点，且离心率为 \(\frac{\sqrt{5}}{2}\) 的双曲线方程.
\end{problem}

\begin{solution}
椭圆的半长轴和半短轴分别为 \(3\) 和 \(2\)，所以其焦半距为 \(c=\sqrt{3^2-2^2}=\sqrt{5}\). 所求双曲线与椭圆有公共焦点，故双曲线的焦点也在 \(x\) 轴上，焦半距仍为 \(\sqrt{5}\). 设双曲线的半实轴为 \(a\)，由离心率条件 \(\frac{c}{a}=\frac{\sqrt{5}}{2}\)，得 \(a=2\). 再由双曲线的关系式 \(c^2=a^2+b^2\)，得 \(b^2=5-4=1\). 因此双曲线方程为 \(\frac{x^2}{4}-y^2=1\).
\end{solution}

\begin{problem}
如图，\(AB\) 是圆 \(O\) 的直径，\(PA\) 垂直于圆 \(O\) 所在的平面，\(C\) 是圆周上不同于 \(A\)，\(B\) 的任一点，求证：平面 \(PAC\) 垂直于平面 \(PBC\).

\bitmapinlinefigure[9.1cm]{img/1986/national_liberal/q17_fig1.png}
\end{problem}

\begin{solution}
因为 \(PA\) 垂直于圆 \(O\) 所在的平面，所以 \(PA\perp AC\) 且 \(PA\perp BC\). 又因为 \(AB\) 是圆的直径，\(C\) 在圆周上，由圆周角定理得 \(\angle ACB=90^{\circ}\)，即 \(AC\perp BC\). 直线 \(PA\) 与 \(AC\) 是平面 \(PAC\) 内相交的两条直线，而 \(BC\) 同时垂直于这两条直线，所以直线 \(BC\) 垂直于平面 \(PAC\). 由于直线 \(BC\) 在平面 \(PBC\) 内，所以平面 \(PBC\) 垂直于平面 \(PAC\)，即平面 \(PAC\perp\) 平面 \(PBC\).
\end{solution}

\begin{problem}
求满足方程 \(|z + 3 - \sqrt{3}\i| = \sqrt{3}\) 的辐角主值最小的复数 \(z\).
\end{problem}

\begin{solution}
令 \(z=x+y\i\)，则原方程表示以 \((-3,\sqrt{3})\) 为圆心、以 \(\sqrt{3}\) 为半径的圆. 圆心到原点的距离为 \(\sqrt{(-3)^2+(\sqrt{3})^2}=2\sqrt{3}\)，圆心对应的辐角为 \(\frac{5\pi}{6}\). 从原点向此圆作切线，设切线与圆心连线的夹角为 \(\delta\)，则在由原点、圆心和切点构成的直角三角形中，\(\sin\delta=\frac{\sqrt{3}}{2\sqrt{3}}=\frac{1}{2}\)，所以 \(\delta=\frac{\pi}{6}\). 因而圆上各点的辐角主值在 \(\frac{5\pi}{6}-\frac{\pi}{6}=\frac{2\pi}{3}\) 与 \(\pi\) 之间，最小值在辐角为 \(\frac{2\pi}{3}\) 的切点处取得.

设该切点对应的复数为 \(z=t\left(\cos\frac{2\pi}{3}+\i\sin\frac{2\pi}{3}\right)\)，其中 \(t>0\). 代入圆的方程，得

\[
\left(3-\frac{t}{2}\right)^2+3\left(\frac{t}{2}-1\right)^2=3,
\]

即 \((t-3)^2=0\)，所以 \(t=3\). 故所求复数为 \(z=3\left(\cos\frac{2\pi}{3}+\i\sin\frac{2\pi}{3}\right)=-\frac{3}{2}+\frac{3\sqrt{3}}{2}\i\).
\end{solution}

\begin{problem}
已知抛物线 \(y^2 = x + 1\)，定点 \(A(3,1)\)，\(B\) 为抛物线上任意一点，点 \(P\) 在线段 \(AB\) 上，且有 \(BP:PA = 1:2\)，当点 \(B\) 在抛物线上变动时，求点 \(P\) 的轨迹方程，并指出这个轨迹为哪种曲线.
\end{problem}

\begin{solution}
设抛物线上的点 \(B\) 为 \(B(t^2-1,t)\)，其中 \(t\in\R\). 由 \(BP:PA=1:2\)，点 \(P\) 将线段 \(BA\) 按内分比 \(1:2\) 分割，所以

\[
P=\frac{2B+A}{3}=\left(\frac{2t^2+1}{3},\frac{2t+1}{3}\right).
\]

由 \(y=\frac{2t+1}{3}\) 得 \(t=\frac{3y-1}{2}\)，再由 \(x=\frac{2t^2+1}{3}\) 得 \(6x=(3y-1)^2+2\). 整理得 \(3y^2-2y-2x+1=0\)，即 \(3\left(y-\frac{1}{3}\right)^2=2\left(x-\frac{1}{3}\right)\). 因此点 \(P\) 的轨迹是顶点为 \(\left(\frac{1}{3},\frac{1}{3}\right)\)、开口向右的抛物线.
\end{solution}

\begin{problem}
甲，乙，丙，丁四个公司承包 \(8\) 项工程，甲公司承包 \(3\) 项，乙公司承包 \(1\) 项，丙，丁两个公司各承包 \(2\) 项，问共有多少种承包方式.
\end{problem}

\begin{solution}
\(8\) 项工程彼此不同，\(4\) 个公司有明确区别. 先从 \(8\) 项工程中选出甲公司承包的 \(3\) 项，再从余下的 \(5\) 项中选出乙公司承包的 \(1\) 项，然后从余下的 \(4\) 项中选出丙公司承包的 \(2\) 项，最后剩下的 \(2\) 项由丁公司承包. 因此承包方式共有

\[
\mathrm C_8^3\mathrm C_5^1\mathrm C_4^2\mathrm C_2^2
=56\times5\times6\times1
=1680
\]

种.
\end{solution}

\begin{problem}
已知 \(\sin A + \sin 3A + \sin 5A = a\)，\(\cos A + \cos 3A + \cos 5A = b\). 求证：

\begin{enumerate}
\item 当 \(b \neq 0\) 时，\(\tan 3A = \frac{a}{b}\)；
\item \(\left(1 + 2\cos 2A\right)^2 = a^2 + b^2\).
\end{enumerate}
\end{problem}

\begin{solution}
\begin{enumerate}
\item 由和差化积公式，

\[
a=\sin A+\sin 3A+\sin 5A=\left(1+2\cos 2A\right)\sin 3A,
\]

\[
b=\cos A+\cos 3A+\cos 5A=\left(1+2\cos 2A\right)\cos 3A.
\]

当 \(b\neq0\) 时，由 \(b=\left(1+2\cos 2A\right)\cos 3A\neq0\)，可知 \(\cos 3A\neq0\)，所以

\[
\tan 3A=\frac{\sin 3A}{\cos 3A}=\frac{\left(1+2\cos 2A\right)\sin 3A}{\left(1+2\cos 2A\right)\cos 3A}=\frac{a}{b}.
\]

\item 由上述两个等式以及 \(\sin^2 3A+\cos^2 3A=1\)，得

\[
a^2+b^2=\left(1+2\cos 2A\right)^2\left(\sin^2 3A+\cos^2 3A\right)=\left(1+2\cos 2A\right)^2.
\]

\end{enumerate}
\end{solution}

\begin{problem}
已知数列 \(\{a_n\}\)，其中 \(a_1 = \frac{4}{3}\)，\(a_2 = \frac{13}{9}\)，且当 \(n \geqslant 3\) 时，\(a_n - a_{n-1} = \frac{1}{3}\left(a_{n-1} - a_{n-2}\right)\).

\begin{enumerate}
\item 求数列 \(\{a_n\}\) 的通项公式；
\item 求 \(\displaystyle\lim_{n\to\infty} a_n\).
\end{enumerate}
\end{problem}

\begin{solution}
\begin{enumerate}
\item 令 \(d_n=a_n-a_{n-1}\). 由已知条件，

\[
d_2=a_2-a_1=\frac{13}{9}-\frac{4}{3}=\frac{1}{9}=\frac{1}{3^2}.
\]

当 \(n\geqslant3\) 时，题设递推关系化为 \(d_n=\frac{1}{3}d_{n-1}\)，所以对 \(n\geqslant2\)，

\[
d_n=\frac{1}{9}\left(\frac{1}{3}\right)^{n-2}=\frac{1}{3^n}.
\]

因此，当 \(n\geqslant2\) 时，

\[
a_n=a_1+\sum_{k=2}^{n}d_k=\frac{4}{3}+\sum_{k=2}^{n}\frac{1}{3^k}=\frac{4}{3}+\frac{1}{6}\left(1-\frac{1}{3^{n-1}}\right)=\frac{3}{2}-\frac{1}{2\cdot3^n}.
\]

当 \(n=1\) 时，右端为 \(\frac{3}{2}-\frac{1}{2\cdot3}=\frac{4}{3}=a_1\)，所以数列的通项公式为 \(a_n=\frac{3}{2}-\frac{1}{2\cdot3^n}\)（\(n\geqslant1\)）.

\item 由通项公式，

\[
\displaystyle\lim_{n\to\infty}a_n=\frac{3}{2}-\displaystyle\lim_{n\to\infty}\frac{1}{2\cdot3^n}=\frac{3}{2}.
\]
\end{enumerate}
\end{solution}
